\section{Results}
\label{sec:results}

Unless stated otherwise, scores are mSCF1 at the matched bin budget
$K_{\text{bin}}$ and use training seed~1. Differences are paired by section.
Practical labels are quoted exactly and the statistics of \cref{sec:stats}
are reported beside them.

\subsection{Baseline reproductions (RQ4)}
\label{sec:res-repro}
\Cref{fig:repro} compares the reproduced baselines with the published values,
each under its own protocol and evaluator.
\begin{itemize}
  \item \textbf{msiPL} reproduced closely on all three collections: GBM 0.414
        against 0.403 published, CAC 0.402 against 0.431, and RCC 0.157 against
        0.159.
  \item \textbf{S3PL on CAC and RCC} was close: 0.592 against 0.601 on CAC at
        seed~1, and 0.41--0.44 against 0.47 on RCC over three seeds.
  \item \textbf{S3PL on GBM was not reproduced:} 0.32 against 0.496 published.
        We used the current public repository with the released settings. The
        paper's supplementary tables were not publicly available, so its exact
        per-dataset settings could not be checked. This gap is analysed in
        \cref{sec:res-s3pl}.
\end{itemize}

\begin{figure}[t]
\centering
\includegraphics[width=\columnwidth]{20_baseline_reproductions.png}
\caption{Reproduced against published collection means for msiPL and S3PL,
each under its own protocol and evaluator. msiPL reproduces on all three
collections; S3PL reproduces on CAC and RCC but not on GBM.}
\label{fig:repro}
\end{figure}

\subsection{Main result: ranking by correlation outperforms IG (RQ1, RQ2)}
\label{sec:res-main}
\Cref{tab:main} gives all rankings scored on the same segmentation and budget.

\textbf{Against legacy msiPL (RQ1).} Both VAE--IG arms beat reproduced msiPL in
every GBM and CAC section: 0.518 and 0.511 against 0.414 on GBM, and 0.528 and
0.560 against 0.402 on CAC. This is statistically supported in all four
comparisons (Holm $p=0.031$; \cref{tab:stats}).
The comparison is of complete pipelines. It does not show which ingredient
causes the gain. Legacy msiPL is also not matched at peak level, so one clause
of the predeclared rule (budget stability) could not be assessed, and the rule
is not called met.

\textbf{IG against posterior correlation (RQ2).} Holding each fitted GMM fixed,
ranking bins by their correlation with the GMM posterior maps outperformed IG
in all 47 valid section--arm comparisons at seed~1 (\cref{fig:igcorr}). The
predeclared rule did not anticipate an outcome this one-sided. Its exact
fall-through wording therefore applies: ``neither predeclared verdict applies,
because IG scored below posterior-$|$PCC$|$ in every section''.

\begin{figure}[t]
\centering
\includegraphics[width=\columnwidth]{05_ig_minus_correlation_by_section.png}
\caption{IG minus posterior correlation (mSCF1) per section, for both VAE arms,
at seed~1. Every valid comparison is negative: correlation with the
segmentation outperforms IG on the same segmentation.}
\label{fig:igcorr}
\end{figure}

\textbf{Across seeds.} Over training seeds 1--3, correlation outperformed IG in
136 of 140 valid bin-level comparisons (\cref{tab:seeds}). The strict per-seed
rule (IG below in at least 7/8 sections) failed only at CAC seed~2, where IG
was below in 6/8 sections, so the predeclared label is ``does not hold
overall''. The direction never reversed at collection level: the mean
difference was negative in all 30 dataset $\times$ arm $\times$ seed $\times$
scale cells.

\textbf{Statistics.} With seeds averaged within section, the difference is
statistically supported in all six dataset--arm comparisons (\cref{tab:stats}).
Correlation scored higher in every section (rank-biserial $r=-1$). The
Hodges--Lehmann differences were $-0.248$ (GBM), $-0.078$ (CAC) and $-0.147$
(RCC) for the centre-only arm. Every raw p-value is the smallest attainable for
its $n$, so this is as strong as eight sections can show. The descriptive mixed
models agree: $-0.24$ [$-0.27$, $-0.22$] on GBM and $-0.14$ [$-0.17$, $-0.12$]
on RCC. They attribute almost no variance to the seed. GBM is negative in every
PDX line. Increasing IG sampling from 12 to 48 or 192 pixels per component
raised scores by at most $+0.024$ on the development sections, which did not
meet the rule for changing the production setting.

\begin{table*}[t]
\caption{Collection mean mSCF1 at $K_{\text{bin}}$, seed~1, bin level. All
rankings except the baselines use the same fitted GMM. Bold: best label-free
ranking per row. The oracle uses the expert masks and is a ceiling.}
\label{tab:main}
\centering\small
\begin{tabular}{@{}llccccccccc@{}}
\toprule
Collection & VAE arm & IG & Posterior corr. & Moran's $I$ & Mean int. & Variance & Legacy msiPL & Random & Oracle \\
\midrule
GBM & centre & 0.518 & \textbf{0.743} & 0.684 & 0.483 & 0.497 & 0.414 & 0.009 & 0.809 \\
GBM & uniform & 0.511 & \textbf{0.744} & 0.684 & 0.483 & 0.497 & 0.414 & 0.009 & 0.809 \\
CAC & centre & 0.528 & 0.606 & 0.605 & \textbf{0.612} & 0.593 & 0.402 & 0.108 & 0.654 \\
CAC & uniform & 0.560 & 0.607 & 0.605 & \textbf{0.612} & 0.593 & 0.402 & 0.108 & 0.654 \\
RCC & centre & 0.412 & 0.568 & \textbf{0.631} & 0.498 & 0.508 & --- & 0.171 & 0.756 \\
RCC & uniform$^{a}$ & 0.410 & 0.596 & \textbf{0.635} & 0.499 & 0.504 & --- & 0.173 & 0.755 \\
\bottomrule
\multicolumn{10}{@{}l}{\footnotesize $^{a}$ 7 of 8 runs: UH0505\_12 failed the predeclared IG completeness check. Legacy msiPL on RCC is not budget-matched (\cref{sec:res-repro}).}
\end{tabular}
\end{table*}

\begin{table}[t]
\caption{IG minus posterior correlation (bin level) across training seeds:
mean $\pm$ SD across seeds 1--3 of the seed-level collection means, and the
number of valid section comparisons in which IG was lower.}
\label{tab:seeds}
\centering\small
\begin{tabular}{@{}llcc@{}}
\toprule
Collection & Arm & IG $-$ corr. & IG lower \\
\midrule
GBM & centre & $-0.244 \pm 0.018$ & 24/24 \\
GBM & uniform & $-0.245 \pm 0.011$ & 23/23 \\
CAC & centre & $-0.077 \pm 0.003$ & 23/24 \\
CAC & uniform & $-0.055 \pm 0.007$ & 22/24 \\
RCC & centre & $-0.144 \pm 0.011$ & 22/23 \\
RCC & uniform & $-0.141 \pm 0.041$ & 22/22 \\
\midrule
\multicolumn{2}{@{}l}{Total} & & 136/140 \\
\bottomrule
\multicolumn{4}{@{}p{0.95\columnwidth}@{}}{\footnotesize Seed-level means
cover the sections valid at that seed. Four seed-arms failed the IG
completeness check. Per-seed tables are in \cref{app:seeds}.}
\end{tabular}
\end{table}

\begin{table}[t]
\caption{Paired statistics for the ranking (F1) and legacy (F2) families, bin
level, sections with seeds averaged. HL: Hodges--Lehmann median difference
with exact 95\% CI. Holm adjusts within each family.}
\label{tab:stats}
\centering\scriptsize
\begin{tabular}{@{}lccccc@{}}
\toprule
Comparison & $n$ & HL [95\% CI] & $r$ & $p$ & Holm $p$ \\
\midrule
\multicolumn{6}{@{}l}{\emph{F1: IG $-$ posterior correlation}} \\
GBM centre & 8 & $-0.248$ [$-0.293$, $-0.186$] & $-1$ & 0.008 & 0.047 \\
GBM uniform & 7 & $-0.243$ [$-0.314$, $-0.159$] & $-1$ & 0.016 & 0.047 \\
CAC centre & 8 & $-0.078$ [$-0.116$, $-0.040$] & $-1$ & 0.008 & 0.047 \\
CAC uniform & 8 & $-0.046$ [$-0.080$, $-0.027$] & $-1$ & 0.008 & 0.047 \\
RCC centre & 7 & $-0.147$ [$-0.207$, $-0.082$] & $-1$ & 0.016 & 0.047 \\
RCC uniform & 6 & $-0.131$ [$-0.204$, $-0.056$] & $-1$ & 0.031 & 0.047 \\
\midrule
\multicolumn{6}{@{}l}{\emph{F2: IG $-$ legacy msiPL (seed 1)}} \\
GBM centre & 8 & $0.102$ [$0.073$, $0.124$] & $1$ & 0.008 & 0.031 \\
GBM uniform & 8 & $0.094$ [$0.057$, $0.136$] & $1$ & 0.008 & 0.031 \\
CAC centre & 8 & $0.127$ [$0.065$, $0.186$] & $1$ & 0.008 & 0.031 \\
CAC uniform & 8 & $0.164$ [$0.097$, $0.208$] & $1$ & 0.008 & 0.031 \\
\bottomrule
\end{tabular}
\end{table}

\textbf{Answer to RQ1 and RQ2.} The pipeline beats legacy msiPL, but the
attribution method is not what carries that gain. Given the same segmentation,
a direct correlation ranking is better than IG on all three datasets. The
useful signal is in the label-free segmentation, not in the gradients.

\subsection{Simple rankings and the metric}
\label{sec:res-simple}
Among the predeclared label-free rankings, the strongest differs by dataset
(\cref{tab:main}):
\begin{itemize}
  \item on GBM, posterior correlation (0.743), with Moran's~$I$ at 0.684;
  \item on CAC, mean intensity (0.612), with correlation, Moran's~$I$ and
        variance all within 0.02;
  \item on RCC, Moran's~$I$ (0.63), above correlation (0.57--0.60).
\end{itemize}
A max-aggregated correlation variant, which was not part of the predeclared
comparison, scores 0.626--0.629 on CAC. It is reported in the appendix only.

CAC is close to saturation at its budget. All methods converge near 0.41 at
twice the budget, including the oracle. This limits how much CAC differences
can reveal. On GBM and RCC, sections whose GMM agrees better with the masks
(balanced accuracy) have correlation scores closer to the oracle (Spearman
$\rho=-0.90$ and $-0.84$ against the gap to the oracle). mSCF1 rewards ion
images that follow the expert regions. A segmentation that resembles those
regions therefore scores well when correlated against, and this alignment is
part of why correlation does well (\cref{sec:discussion}).

\subsection{Spatial context in the encoder and the segmentation (RQ3)}
\label{sec:res-spatial}
\begin{itemize}
  \item \textbf{Uniform $3\times3$ context.} Under IG, context ``helps'' on
        CAC: $+0.023$ (bin) and $+0.027$ (P1) averaged over three seeds, positive
        in 8/8 sections. It is not statistically supported (Holm $p=0.13$), and
        it is 0.000 under correlation ranking. On GBM the three-seed result is
        ``incomplete (7/8)'', with a mean of $+0.000$ (\cref{fig:gbmpdx}).
  \item \textbf{Spatial specificity.} If the CAC gain came from correct local
        placement, real neighbours should beat shuffled nonlocal ones. They
        tied: 0.542 for real $3\times3$ against 0.543 for shuffled. On GBM,
        zeroed context scored highest (0.459, against 0.446 for real).
  \item \textbf{Larger windows.} $7\times7$ and $9\times9$ windows did not meet
        the ``helps'' rule. Both were $0.035$ below $3\times3$ on GBM (positive
        in 2/8, ``mixed'') and ``equivalent'' on CAC (\cref{fig:variants}).
  \item \textbf{Learned attention} was ``equivalent'' to uniform averaging on
        both collections ($+0.011$ GBM, $-0.006$ CAC). With real neighbours its
        weights stayed almost uniform (normalised entropy 0.98--0.999). With
        shuffled neighbours it learned selective weights, so the mechanism
        works; real neighbours are simply near-copies of the centre (median
        cosine about 0.98). Masking the centre during training did not make
        attention selective, and it lowered GBM scores by $0.039$.
  \item \textbf{Spatially smoothed segmentation (Test~S)} was never labelled
        ``helps''. Four comparisons were ``equivalent'', one ``does not meet
        the `helps' rule (labelled mixed)'' (GBM uniform, $+0.029$, driven by
        two sections), and RCC uniform was ``incomplete'' (\cref{app:extra}).
\end{itemize}

\begin{figure}[t]
\centering
\includegraphics[width=\columnwidth]{17_all_neighbourhood_variants.png}
\caption{All neighbourhood variants relative to uniform $3\times3$ context
under IG. No variant gives a consistent, spatially specific gain.}
\label{fig:variants}
\end{figure}

\begin{figure}[t]
\centering
\includegraphics[width=\columnwidth]{08_gbm_context_by_patient.png}
\caption{Uniform minus centre-only IG score on GBM, grouped by PDX line. The
direction differs between lines.}
\label{fig:gbmpdx}
\end{figure}

\textbf{Answer to RQ3.} Neighbourhood information entering the encoder or the
segmentation gave no consistent, spatially specific gain in this pipeline. The
spatial information that did help sits in the ranking statistic: correlation
with spatially coherent segment maps, and Moran's~$I$, which ranks bins by how
spatially structured their ion images are.

\subsection{Is the nonlinear encoder needed? (RQ2)}
\label{sec:res-encoder}
Replacing the VAE latent space with a 5-component PCA, while keeping the GMM and
correlation ranking (Test~V), gave the following VAE $-$ PCA-5 differences:
\begin{itemize}
  \item GBM: $+0.023$, higher in 5/8 sections, ``mixed'';
  \item CAC: $-0.013$, ``mixed'';
  \item RCC: ``equivalent''.
\end{itemize}
Averaging the VAE over three seeds did not change the verdicts
(\cref{fig:vaepca}). VAE segmentations matched the masks slightly better
(balanced accuracy 0.868 against 0.839 on GBM), without reliably better peaks.
Reconstruction did not separate the two either. PCA-5 had \emph{lower}
reconstruction error than the VAE on all eight GBM sections, and slightly
higher error on 7/8 CAC sections. This matches the finding of
\citet{mazarura2025autoencoder} that reconstruction quality and downstream
usefulness can rank models differently.
The number of GMM components was not validated without labels. BIC selected
the largest candidate ($K=6$) in every tested case, with no interior minimum,
and a finer segmentation lowered GBM scores by about 0.30.

\begin{figure}[t]
\centering
\includegraphics[width=\columnwidth]{vae_vs_pca.png}
\caption{VAE minus PCA-5 posterior-correlation score per section, using the
same GMM and ranking. The verdicts are ``mixed'' (GBM, CAC) and
``equivalent'' (RCC).}
\label{fig:vaepca}
\end{figure}

\subsection{S3PL: neighbour dependence, reproduction gap and seeds (RQ3, RQ4)}
\label{sec:res-s3pl}
\textbf{Neighbour dependence.} Retraining S3PL with each patch filled by copies
of the centre spectrum collapsed it from 0.591 to 0.18 and 0.15 on CAC (two
seeds, 8/8 sections). Retraining with shuffled nonlocal neighbours recovered
most of the score (0.55 against 0.59). Correct placement added about $+0.04$
inconsistently (4/8 and 3/8 sections at least $+0.02$), so the label is
``mixed'' (Holm $p=0.63$). S3PL needs neighbour input under its patch
objective, but a consistent benefit from \emph{spatially correct} neighbours
was not established. Training S3PL for longer made it worse: from epoch 10 to
100, CAC fell from 0.591 to 0.335 and GBM from 0.316 to 0.212. Epoch 10 was
best in 15 of 16 sections, so the released 10-epoch protocol does not
handicap it.

\textbf{Reproduction gap.} With ten released seeds per section
(\cref{fig:s3plseeds}), S3PL scored $0.319 \pm 0.042$ on GBM and
$0.547 \pm 0.038$ on CAC (mean $\pm$ SD across seeds of the collection means).
\begin{itemize}
  \item \textbf{GBM:} the published 0.496 is about four SD above our mean and
        above all ten seed-level means (maximum 0.374). It lies above every one
        of our seeds in 6/8 sections. Taking the best of ten seeds in every
        section reaches only 0.434. Labels: ``outside our seed distribution'';
        ``not explained by seed selection alone''.
  \item \textbf{CAC:} the published 0.601 is just outside our range (maximum
        0.591). Best-of-ten reaches 0.5998, which is within the published
        two-decimal rounding. So CAC is consistent with a favourable seed or
        best-of-several reporting, but this is not shown.
  \item \textbf{Preprocessing:} replacing the released spatial-maximum
        normalisation with the TIC normalisation described in the paper did not
        raise the GBM mean (0.303 against 0.324; HL $-0.012$ [$-0.154$,
        $0.099$]). Label: ``does not explain''. The seed-1 TIC run predates this
        test.
\end{itemize}
The remaining candidate explanations are untested: a different data
preparation for specific sections, or hyperparameters tuned on the evaluation
sections, which the paper describes.

\textbf{Why seeds change S3PL.} About half of the score variance is between
sections, about 0--1\% is a seed main effect, and about 47\% is a section
$\times$ seed interaction. Seeds do not shift whole collections; they switch
particular sections between distinct outcomes:
\begin{itemize}
  \item GBM108\_positive moves between 0.02 and 0.46;
  \item CAC 520TopL moves between 0.17 and 0.71;
  \item in flipping sections, peak lists from different seeds can overlap by as
        little as 1\%.
\end{itemize}
Separating the seed into its parts on three such sections showed the
following:
\begin{itemize}
  \item weight initialisation dominated data order (SD 0.222 against 0.001 on
        GBM108\_positive, and 0.128 against 0.010 on GBM12\_2), giving the label
        ``initialisation dominant'';
  \item on GBM108\_negative, data order also mattered (ratio 1.55);
  \item same-seed reruns were bit-identical, so GPU nondeterminism played no
        part (``not dominated by nondeterminism'').
\end{itemize}
Under TIC normalisation the per-section seed SD fell from 0.061 to 0.011 and
the switching disappeared. Released S3PL therefore reaches at least two
distinct solutions on these sections, and its starting weights mostly decide
which one.

\begin{figure*}[t]
\centering
\includegraphics[width=0.9\textwidth]{s3pl_released_seed_distributions.pdf}
\caption{S3PL mSCF1 per section over ten released seeds (dots), seed mean
(bar) and published value (diamond), using S3PL's own evaluator. GBM: the
published value is above every seed in 6/8 sections (``outside our seed
distribution''). CAC: just outside, within the published rounding. One CAC run
(240TopL, seed~7) returned 256 of 257 peaks, as the released code does when
fewer distinct bins are selected.}
\label{fig:s3plseeds}
\end{figure*}

\textbf{Answer to RQ4.} msiPL reproduces. S3PL reproduces on CAC and RCC but
not on GBM, and neither seed selection nor the documented preprocessing
explains the GBM gap. S3PL is highly seed-sensitive in specific sections
because its initialisation selects between distinct solutions, so
single-seed comparisons against it are unreliable.

\subsection{RCC: held-out confirmation and its qualifier}
\label{sec:res-rcc}
On the held-out RCC collection, correlation outperformed IG in every valid run
(centre $-0.156$, 0/8 positive; uniform $-0.185$, 0/7), and Moran's~$I$ was
the strongest label-free ranking (0.63--0.64). This confirms the main result
on a dataset not used for development. One qualifier is necessary. The
unannotated class covers 45--78\% of pixels, and when only cancer and normal
pixels are scored, the advantage shrinks:
\begin{itemize}
  \item for the centre-only arm, correlation beats IG in 5/8 runs and falls
        below mean intensity, variance and Moran's~$I$;
  \item for the uniform arm, correlation beats IG in 6/7 runs.
\end{itemize}
The three-class RCC result is therefore partly driven by separating annotated
tissue from the unannotated region.

\subsection{Runtime}
\label{sec:res-runtime}
\Cref{tab:runtime} reports training and ranking times. Once a segmentation
exists, correlation takes about a second or less per section, while IG takes
seconds to minutes. Every timed ranking reproduced the saved ranking exactly.
Cached 100-epoch VAE training takes longer than S3PL's 10-epoch protocol, but
remains practical. The proposal asked for memory and runtime against S3PL. On
CAC section 160TopL, three complete native-workflow repeats (training,
selection and scoring) took 122--135\,s for the VAE against 151--174\,s for
S3PL, with peak memory of 1.15--1.18\,GiB against 1.31--1.32\,GiB (GPU
237 against 315\,MiB). These are measurements under each method's native protocol,
not an equal-work efficiency comparison.

\begin{table}[t]
\caption{Median runtime per section. Training: cached 100-epoch VAE (seeds
2--3) and 10-epoch S3PL, on GPU. Ranking: CPU, two threads, given a trained
model; IG includes the GMM fit and faithfulness checks.}
\label{tab:runtime}
\centering\small
\begin{tabular}{@{}lccc@{}}
\toprule
 & GBM & CAC & RCC \\
\midrule
VAE training (min) & 32--35 & 1.3--1.4 & 2.2 \\
S3PL training (min) & 13.7 & 0.9 & 1.6 \\
\midrule
Posterior corr.\ (s) & 1.12 & 0.028 & 0.11 \\
Moran's $I$ (s) & 8.4 & 0.18 & 0.75 \\
Mean intensity (s) & 0.10 & 0.003 & 0.013 \\
IG (s) & 770 & 6.0 & 372 \\
\bottomrule
\multicolumn{4}{@{}p{0.95\columnwidth}@{}}{\footnotesize Ranking CPUs: i5-14400
(GBM, CAC) and i9-10940X (RCC, shared node), each matching the CPU that produced
the saved rankings.}
\end{tabular}
\end{table}
